\documentclass{course_sheet}
\usepackage{amsmath,amssymb}
\usepackage{longtable}
\usepackage{xcolor}
\setlength{\parindent}{0pt}
\setlength{\emergencystretch}{2em}
\begin{document}

\title{Exercise: Diagnosing an Infection with Temperature}
\author{}
\date{}
\maketitle
\subsection*{Learning objectives}
By the end of this exercise, you should be able to:

\begin{itemize}
\item
Calculate and interpret the \textbf{\textit{Gini index}} for a data set.

\item
Understand how changing a \textbf{\textit{classification threshold}} affects the separation between two signal and background.

\item
Calculate \textbf{\textit{sensitivity}} (true-positive rate).

\item
Calculate \textbf{\textit{signal efficiency}} and \textbf{\textit{background efficiency}}.

\item
Understand the trade-off between signal efficiency and background rejection.

\end{itemize}
\subsection*{Scenario}
A hospital is testing whether a patient's \textbf{\textit{body temperature}} can be used to diagnose an infection.

Ten patients have been examined. After a more definitive medical test, we know whether each patient actually had an infection. We will call:

\begin{itemize}
\item
\textbf{\textit{Signal:}} patients who actually have an infection.

\item
\textbf{\textit{Background:}} patients who do not have an infection.

\end{itemize}
The hospital considers a patient \textbf{\textit{positive for infection}} if their temperature is \textbf{\textit{above a chosen threshold}}.

\subsubsection*{Patient data}
\begingroup
\small
\setlength{\tabcolsep}{4pt}
\begin{longtable}{l|l|l}
Patient & Temperature (°C) & Actually infected? \\[1.5pt]
\hline
\endfirsthead
Patient & Temperature (°C) & Actually infected? \\[1.5pt]
\hline
\endhead
1 & 36.5 & No \\[1.5pt]
2 & 36.8 & No \\[1.5pt]
3 & 37.0 & No \\[1.5pt]
4 & 37.2 & Yes \\[1.5pt]
5 & 37.5 & No \\[1.5pt]
6 & 37.7 & Yes \\[1.5pt]
7 & 38.0 & Yes \\[1.5pt]
8 & 38.2 & Yes \\[1.5pt]
9 & 38.5 & No \\[1.5pt]
10 & 39.0 & Yes \\[1.5pt]
\end{longtable}
\endgroup
There are therefore \textbf{\textit{5 infected patients (signal)}} and \textbf{\textit{5 non-infected patients (background)}}.

\clearpage
\section*{Part I - Choosing a temperature threshold}
Suppose the hospital initially chooses a threshold of

\[
T > 37.5^\circ\mathrm{C}.
\]

A patient is classified as \textbf{\textit{positive}} when their temperature is above this value.

\subsubsection*{Exercise 1.1}
\begin{enumerate}
\item
Which patients are classified as positive?

\item
Among these patients, how many are:

\begin{itemize}
\item
\textbf{\textit{True positives (TP)}} - infected and classified as positive?

\item
\textbf{\textit{False positives (FP)}} - not infected but classified as positive?

\end{itemize}
\item
How many patients are classified as negative?

\item
Among the negative patients, how many are:

\begin{itemize}
\item
\textbf{\textit{True negatives (TN)}}?

\item
\textbf{\textit{False negatives (FN)}}?

\end{itemize}
\end{enumerate}
Complete the confusion matrix:

\begingroup
\small
\setlength{\tabcolsep}{4pt}
\begin{longtable}{l|l|l}
\rule{0pt}{2.4ex} & Actually infected & Not infected \\[3pt]
\hline
\endfirsthead
\rule{0pt}{2.4ex} & Actually infected & Not infected \\[3pt]
\hline
\endhead
\textbf{\textit{Classified positive}} & TP = \_\_\_ & FP = \_\_\_ \\[1.5pt]
\textbf{\textit{Classified negative}} & FN = \_\_\_ & TN = \_\_\_ \\[1.5pt]
\end{longtable}
\endgroup
\clearpage
\section*{Part II - Sensitivity}
\textbf{\textit{Sensitivity}} measures the fraction of all truly infected patients that the diagnostic method successfully identifies.

The formula is

\[
{
\text{Sensitivity}
=
\frac{\text{TP}}{\text{TP}+\text{FN}}
}
\]

\subsubsection*{Exercise 2.1}
Using the threshold $T>37.5^\circ\mathrm{C}$:

\begin{enumerate}
\item
Calculate the sensitivity.

\item
Express your answer as a percentage.

\item
What does this percentage mean in the context of the hospital?

\end{enumerate}
\clearpage
\section*{Part III - Signal and background efficiency}
In a classification problem, we can interpret the infected patients as \textbf{\textit{signal}} and the non-infected patients as \textbf{\textit{background}}.

\subsubsection*{Signal efficiency}
The \textbf{\textit{signal efficiency}} is the fraction of all signal events that pass the selection:

\[
{
\epsilon_S =
\frac{\text{number of signal events passing}}
{\text{total number of signal events}}
}
\]

Notice that, in this example,

\[
\epsilon_S = \text{sensitivity}.
\]

\subsubsection*{Background efficiency}
The \textbf{\textit{background efficiency}} is the fraction of all background events that pass the selection:

\[
{
\epsilon_B =
\frac{\text{number of background events passing}}
{\text{total number of background events}}
}
\]

\subsubsection*{Exercise 3.1}
For the threshold $T>37.5^\circ\mathrm{C}$:

\begin{enumerate}
\item
Calculate the signal efficiency.

\item
Calculate the background efficiency.

\item
What fraction of the background is rejected?

\item
Is it possible to increase the signal efficiency while decreasing the background efficiency? Investigate this by changing the temperature threshold.

\end{enumerate}
\clearpage
\section*{Part IV - The Gini index}
The \textbf{\textit{Gini index}} is a measure of how mixed two classes are within a sample.

For two classes, it is defined as

\[
{
G = 1-p_S^2-p_B^2
}
\]

where:

\begin{itemize}
\item
$p_S$ is the fraction of patients who are signal,

\item
$p_B$ is the fraction of patients who are background,

\item
$p_S+p_B=1$.

\end{itemize}
A sample containing only one class has

\[
G=0,
\]

while a sample containing equal amounts of signal and background has

\[
G=0.5.
\]

\subsubsection*{Exercise 4.1}
Consider \textbf{\textit{all 10 patients together}}.

\begin{enumerate}
\item
What is $p_S$?

\item
What is $p_B$?

\item
Calculate the Gini index.

\end{enumerate}
Now consider the patients classified as \textbf{\textit{positive}} using $T>37.5^\circ\mathrm{C}$.

\begin{enumerate}
\setcounter{enumi}{3}
\item
How many signal patients are in this group?

\item
How many background patients are in this group?

\item
Calculate $p_S$ and $p_B$ for this group.

\item
Calculate its Gini index.

\end{enumerate}
Finally, consider the patients classified as \textbf{\textit{negative}}.

\begin{enumerate}
\setcounter{enumi}{7}
\item
Calculate the Gini index for the negative group.

\item
Compare the Gini indices of the positive and negative groups with the Gini index of the original sample.

\item
What does a lower Gini index tell you about the composition of a group?

\end{enumerate}
\clearpage
\section*{Part V - Investigating the threshold}
The choice of $37.5^\circ\mathrm{C}$ was arbitrary. Let's see what happens when we change it.

Calculate the signal and background efficiencies for several thresholds.

\subsubsection*{Exercise 5.1}
\begin{enumerate}
\item
What happens to the \textbf{\textit{signal efficiency}} as the threshold is increased?

\item
What happens to the \textbf{\textit{background efficiency}}?

\item
Why is there a trade-off between keeping signal and rejecting background?

\item
Which threshold gives the largest signal efficiency?

\item
Which threshold gives the smallest background efficiency?

\item
Is there a single threshold that simultaneously maximizes signal efficiency and minimizes background efficiency? Explain.

\end{enumerate}
\clearpage
\section*{Part VI - Challenge: Finding a useful separation}
Suppose the hospital wants a diagnostic method that keeps as many infected patients as possible while rejecting as many non-infected patients as possible.

One way of visualizing the performance is to plot signal efficiency against background efficiency.

\subsubsection*{Exercise 6.1}
Using the results from Part 5:

\begin{enumerate}
\item
Make a scatter plot with \textbf{\textit{background efficiency on the x-axis}} and \textbf{\textit{signal efficiency on the y-axis}}.

\item
Put a label next to each point showing its temperature threshold.

\item
Describe what happens as the temperature threshold is increased.

\item
Explain why the choice of threshold depends on the consequences of false positives versus false negatives.

\end{enumerate}
\clearpage
\section*{Part VII - Short conceptual questions}
\subsubsection*{Exercise 7.1}
Consider the following questions:

\begin{itemize}
\item
If missing an infected patient is considered very serious, would the hospital tend to prefer a \textbf{\textit{higher or lower}} temperature threshold?

\item
What would happen to the number of false positives if the threshold were lowered?

\item
What would happen to the sensitivity?

\item
Why might a hospital not simply choose the threshold that gives the smallest background efficiency?

\item
How does this example illustrate the general problem of \textbf{\textit{classification}} in machine learning and experimental physics?

\end{itemize}
\subsubsection*{Key concepts}
\begingroup
\small
\setlength{\tabcolsep}{4pt}
\begin{longtable}{l|l}
Concept & Meaning in this exercise \\[1.5pt]
\hline
\endfirsthead
Concept & Meaning in this exercise \\[1.5pt]
\hline
\endhead
\textbf{\textit{Signal}} & Patients who actually have an infection \\[1.5pt]
\textbf{\textit{Background}} & Patients who do not have an infection \\[1.5pt]
\textbf{\textit{Threshold}} & Temperature used to classify a patient as positive \\[1.5pt]
\textbf{\textit{TP}} & Infected patient correctly classified as positive \\[1.5pt]
\textbf{\textit{FP}} & Non-infected patient incorrectly classified as positive \\[1.5pt]
\textbf{\textit{FN}} & Infected patient incorrectly classified as negative \\[1.5pt]
\textbf{\textit{TN}} & Non-infected patient correctly classified as negative \\[1.5pt]
\textbf{\textit{Sensitivity}} & Fraction of infected patients identified \\[1.5pt]
\textbf{\textit{Signal efficiency}} & Fraction of signal passing the selection \\[1.5pt]
\textbf{\textit{Background efficiency}} & Fraction of background passing the selection \\[1.5pt]
\textbf{\textit{Gini index}} & Measure of how mixed the two classes are \\[1.5pt]
\end{longtable}
\endgroup

\end{document}
